\section{Introduction}\label{sec:introduction}

A group has type $F_n$ if it admits a classifying CW complex with finite
$n$-skeleton, and type $\Finfty$ if it admits one classifying complex
with finitely many cells in every dimension. Here a classifying complex
is a connected CW complex with the given fundamental group and
contractible universal cover; it may have infinitely many dimensions.
Type $F_2$ is equivalent to finite presentability.

\begin{theorem}\label{thm:main}
There exists a group $H$ of type $\Finfty$ such that every finitely
presented group admits an injective homomorphism into $H$.
\end{theorem}

The common target gives an exact description of its finitely generated
subgroups. A finitely generated group is recursively presented if it
admits a presentation on finitely many generators with a recursively
enumerable set of relators.

\begin{corollary}\label{cor:recursive}
A finitely generated group embeds in $H$ if and only if it is
recursively presented.
\end{corollary}

\begin{proof}
If $G$ is finitely generated and recursively presented, Higman's
embedding theorem places it in a finitely presented
group~\cite{Higman1961}, which embeds in $H$ by
Theorem~\ref{thm:main}. Conversely, $H$ is finitely presented because
it has type $\Finfty$. For a finitely generated subgroup $G\le H$,
represent its generators by words in a finite presentation of $H$.
Enumerate all words in the chosen generators, and dovetail finite
derivations from the ambient relators proving that their substituted
words are trivial. Recording the successful words enumerates exactly
the defining relators of $G$.
\end{proof}

This answers the $\Finfty$ form of the higher-dimensional Higman
embedding question, stated for finitely generated recursively presented
groups in~\cite[Question~1.3]{FFZ2026}. The theorem imposes no
word-problem hypothesis on the input.

\subsection*{Context and construction}

For a commutative unital coefficient ring $\Lambda$, a group $G$ has
type $FP_n(\Lambda)$ if the trivial $\Lambda G$-module $\Lambda$ has a
projective resolution with finitely generated terms through degree $n$;
write $FP_n$ for $\Lambda=\Z$. For $n=\infty$, finite generation is
required in every degree. Leary proved that every countable group embeds
in a group of type $FP_2$~\cite{Leary2018}. Fournier-Facio and Zaremsky
show that, for each $n\in\N\cup\{\infty\}$, embeddings of all finitely
generated recursively presented groups into recursively presented
groups of type $FP_n$ would imply embeddings of all finitely presented
groups into groups of type $F_n$~\cite[Theorem~A]{FFZ2026}.
They also prove that the rope-trick groups of their Definition~1.5 fail
type $FP_3(\mathbb Q)$ for infinite inputs with nontrivial
free-presentation kernel~\cite[Definition~1.5 and Theorem~B]{FFZ2026}.
This obstruction is specific to those rope-trick groups.

The proof of their conditional theorem uses a universal finitely
presented group, an embedding through an intermediate overgroup and
back, and an ascending HNN extension. They relate this strategy to
earlier work of Baumslag--Dyer--Miller and
Fournier-Facio--L\"oh--Moraschini~\cite[Section~2]{FFZ2026}.

For the universal finitely presented group $U$ of
Lemma~\ref{lem:universal}, Proposition~\ref{prop:finite-envelope}
gives a finite presentation of a group $\widehat W$ equipped with two
finite families of homomorphisms from $U$, sharing a distinguished map
$\iota:U\to\widehat W$. Faithful models in tensor product algebras
verify that $\iota$ is injective. Finite presentability of $\widehat W$ then lets
universality supply an embedding $\kappa:\widehat W\to U$:
\[
\begin{tikzcd}[column sep=large]
 U \arrow[r,hook,"\iota"]
   \arrow[rr,bend right=24,hook,"f=\kappa\iota"']
 & \widehat W \arrow[r,hook,"\kappa"] & U.
\end{tikzcd}
\]
Postcomposing the two families by $\kappa$ gives homomorphisms from $U$
to itself, with the injective endomorphism $f$ as their common
distinguished map (Corollary~\ref{prop:self-map}). The group $U$ embeds
in the associated ascending HNN extension
\[
 H=T(U,f)=\langle U,t\mid t^{-1}ut=f(u)\ (u\in U)\rangle,
\]
which is finitely presented by Lemma~\ref{lem:torus-presentation}.
Theorem~\ref{thm:torus} proves that $H$ has type $\Finfty$.

The finite idempotent systems, the group enlargement, and the
ascending-torus criterion are shared with the
companion~\cite[Sections~2.1, 2.3--2.5, and~3.1]{OpenAIOvergroups}.
We give local proofs of the algebraic identities and geometric
criterion.

The geometric step begins with a uniform filling statement for ordinary
integral chains. In the homogeneous chain model, restricting relative
vertex differences to a finite subset of $U$ leaves only finitely many
simplex types up to common left translation in each degree. For each
positive degree and finite input set of relative differences,
Proposition~\ref{prop:uniform-fillings} chooses one exponent and one
finite output set. That iterate of $f$ sends every integral cycle using
the input differences to the boundary of a chain using the output
differences. The filling lengths and coefficients may vary with the cycle.

The finite diagrams supply these fillings by induction. For each
$\ell\in\{2,3\}$, the diagram is indexed by idempotents over
$\F_\ell$. Orthogonal sums correspond to pointwise products of
homomorphisms with elementwise commuting images.
Lemma~\ref{lem:idempotents} says that an additive assignment has its
value at the identity annihilated by $\ell$. In degrees $n\ge2$, the
chain diagonal includes mixed terms from lower degrees, so the
corresponding cycle classes are not automatically additive. Using
lower-degree fillings, the induction constructs, below degree $n$, a
controlled chain homotopy from the map induced by an iterate of $f$ to
the constant map. Lemma~\ref{lem:controlled-additivity} uses it to
remove the mixed terms. For each input cycle, both diagrams then give
additive assignments into the homology of the same bounded chain
complex. Their common value at the identity idempotent is represented
by the image of the input cycle under an iterate of $f$. It is
annihilated by both $2$ and $3$, and hence vanishes.

For an arbitrary set $\mathcal I$, put $M=(\Z H)^{\mathcal I}$. A cycle
in the tensor complex computing $H_i(U,M)$, $i>0$, decomposes in each
coordinate into finitely many ordinary cycles in left $U$-cosets, with
one finite set of simplex types common to all coordinates. The uniform
filling statement supplies one exponent and output set for these cycles,
so their fillings reassemble into one tensor chain. In the ascending HNN
sequence of Section~\ref{sec:products}, the map $F$ combines $f$ on
chains with multiplication by the stable letter on coefficients. The
fillings make $F$ locally nilpotent on $H_i(U,M)$ for $i>0$: every class
is killed by some power of $F$. Thus $1-F$ is invertible there, and
exactness gives $H_n(H,M)=0$ for every $n\ge2$. Since $H$ is finitely
presented, Lemma~\ref{lem:product-criterion} turns this vanishing into
geometric finiteness. At each stage, the cellular cycle module is
finitely generated; via Hurewicz, those generators have sphere
representatives. Attaching cells along these spheres stage by stage
produces one classifying complex with finitely many cells in every
dimension.

\subsection*{Conventions}

All groups are discrete. Group homology has integral coefficients
unless stated otherwise, and coefficient modules are right modules.
Products indexed by arbitrary sets always mean direct products, not
direct sums.
